\documentclass[%
	type=document,%
  	style=article,% book, article
  	media=print,
  	pages=oneside,%
  	prefixLecturer=Dozenten:,
  	%draft,
]{unihildesheim}
\usepackage[utf8]{inputenc}
\usepackage[ngerman]{babel}
\graphicspath{{../img/}}

% META
\Contact{uni@jens-bertram.net}
\Copyright{cc:by-nc-sa}
\Date{08.07.2012}{08. Juli 2012}
\Lecturer{Gertrud Faaß PhD,\\Dr. Ralph Kölle}
\Seminar{Formalisierung\\[0.5em]{\small Fachbereich III -- Sprach- und
Informationswissenschaften \\Institut für Informationswissenschaft und
Sprachtechnologie}}
\Semester{Sommersemester 2012}
\Title{Hausaufgabe 6}
\Subtitle{Erstellen von Bi- und Trigrammen und
frequenzbasierte Untersuchungen am Beispiel von Zeitungstexten zum Thema
Ministerpräsident}
\Author{J. Bertram, S. Kastner}{Jens Bertram {\small(214508)}\\Sebastian Kastner
{\small(202322)}}
% /META

\begin{document}

\Titlepage
\cleardoublepage
\TOC
\cleardoublepage

%% Wir binden das Uni-Logo auf der Titelseite ein
\begin{titlepage}
\begin{center}
\begin{figure}[htbp]
\begin{center}
\includegraphics[width=3.5cm]{img/UniLogo.png}
\end{center}
\end{figure}

% Name es Instituts
\vspace*{\fill}{
	Stiftung Universität Hildesheim \\ 
	Fachbereich III -- Sprach- und Informationswissenschaften
	\\Institut für Informationswissenschaft und Sprachtechnologie}

\vfill {{\Large Hausaufgabe 6 \\[5mm] Erstellen von Bi- und Trigrammen und
frequenzbasierte Untersuchungen am Beispiel von Zeitungstexten zum Thema
Ministerpräsident}}

{\normalsize \textbf {Sebastian Kastner, 202322, Email:
sebastian\_kastner@gmx.net}}

\vfill{Studiengang:\\ International Information Engineering}
\vfill {Dozenten: Gertrud Faaß, Ralph Koelle\\

(Institut für Informationswissenschaft und Sprachtechnologie)\\}
\end{center}
\end{titlepage}

\thispagestyle{empty}

% CQP:
% http://bulba.sdsu.edu/cqpman.ps
% http://www.ims.uni-stuttgart.de/projekte/CorpusWorkbench/CQPUserManual/HTML/
% http://www.ims.uni-stuttgart.de/projekte/CorpusWorkbench/CQPUserManual/HTML/node27.html#SECTION00830000000000000000
% http://www.ims.uni-stuttgart.de/projekte/CorpusWorkbench/CQPTutorial/cqp-tutorial.pdf

% STTS:
% http://www.ims.uni-stuttgart.de/projekte/corplex/TagSets/stts-table.html

% Adjektiv+Nomen ohne McAllister (Top 10) als LaTeX:
%	// Korpus öffnen
%	MPRAES;
%	// LaTeX Ausgabe Mode
%	set PrintMode latex;
%	// Query
%	@[pos='ADJ(A|D)']@[pos='N(E|N)' & word!='McAllister'];
%	// Filter, Ausgabe in Datei tabelle1.tex
%	group Last matchend word by match lemma > "tabelle1.tex";

% Häufigste Nomina ohne McAllister (Top 10) als LaTeX
%	// Korpus öffnen
%	MPRAES;
%	// LaTeX Ausgabe Mode
%	set PrintMode latex;
%	// Query
%	@[pos='N(E|N)' & word!='McAllister'];
%	// Filter, Ausgabe in Datei tabelle2.tex
%	group Last target lemma > "tabelle2.tex";

% Adjektive mit Suchwort (McAllister) (Top 10) als LaTeX
%	// Korpus öffnen
%	MPRAES;
%	// LaTeX Ausgabe Mode
%	set PrintMode latex;
%	// Query
%	@[pos='ADJ(A|D)']@[word='McAllister'];
%	// Filter, Ausgabe in Datei tabelle3.tex
%	group Last match lemma > "tabelle3.tex";

\section{CQP Ausgaben}

Vor der Ausführung der hier ausgeführten Befehle wurde mit dem Kommando
\it{set PrintMode latex;} das Ausgabeformat auf \LaTeX umgestellt, damit die
Ergebnisse ohne weitere Umformungen weiterverwendet werden konnten.


\subsection{Häufigste Adjektiv-Nomen Paare}
Für die Bestimmung der häufigsten Adjektiv-Nomen Paare, die nicht den
ursprünglichen Suchbegriff "`Ministerpräsident"' enthalten, wurden folgende
Befehle verwendet:
\begin{verbatim}
	> adj_n = @[pos='ADJ(A|D)']@[pos='N(E|N)' & word!="Ministerpräsident.*"];
	> group adj_n matchend lemma by match lemma cut 4 > "tabelle1.tex";
\end{verbatim}

Mit dem ersten Befehl werden alle Adjektiv-Nomen Paare aus dem Korpus ausgelesen
und in der Variable adj\_n gespeichert. Im Folgenden soll der ausgeführte Befehl
kurz erläutert werden:
\begin{enumerate}
  \item \it{@[pos='ADJ(A|D)']} \\
			Das erste Wort muss mit dem POS Tag ADJA oder ADJD getaggt sein
  \item \it{@[pos='N(E|N)' \& word!="Ministerpräsident.*"]]} \\
			Das darauffolgende Wort muss mit dem POS Tag NE oder NE getaggt sein und darf
			nicht gleich Ministerpräsident sein. Durch die Zeichen "`.*"' am Ende von
			Ministerpräsident werden auch deklinierte Formen, wie zum Beispiel
			Ministerpräsidenten ausgeschlossen.
\end{enumerate}

Mit dem zweiten Befehl wurden die Ergebnisse des ersten Befehls nach der
Lemmakombination des ersten (also des Adjektivs) und des zweiten gematchten
Wortes (also das Nomen) gruppiert und nach der Häufigkeit der jeweiligen
Kombination absteigend angeordnet. Da nur die Ausgabe der 10 häufigsten
Kombinationen gefordert war, wurde die Anzahl der ausgegeben Kombinationen
mit \it{cut} beschränkt. Durch \it{cut 4} wurde die Ausgabe auf Kombinationen
beschränkt, die mindestens 4 mal auftreten. Der Wert 4 wurde zuvor manuell
bestimmt, indem der obige Befehl ohne den \it{cut} Befehl verwendet wurde. Die
Ausgabe wurde in die Datei tabelle1.tex umgeleitet.

Das Ergebnis ist der Tabelle \ref{tab:adjektive_nomina} zu entnehmen.

% tabelle 1: adjektive+nomen top10
\begin{table}[htpb]\label{t}
	\center
	\begin{tabularx}{0.6\textwidth}{llr}
		\toprule
		\textbf{Adjektiv} & \textbf{Nomen} & \textbf{Quantität}\\
		\midrule
		% tabelle muss auf nur-Zeilen gekürzt werden!
		%\begin{tabular}{llr}
vergangen
 & Woche
 & 14 \\
 & Jahr
 & 8 \\
öffentlich-rechtlich
 & Rundfunk
 & 7 \\
norddeutsch
 & Landesbank
 & 6 \\
grün
 & Gentechnik
 & 5 \\
rein
 & Spekulation
 & 5 \\
liberaldemokratisch
 & Partei
 & 4 \\
klug
 & Entscheidung
 & 4 \\
offen
 & Denkmal
 & 4 \\
früh
 & Außenminister
 & 4 \\
% dritt
%  & Wahlgang
%  & 4 \\
% neu
%  & Variante
%  & 4 \\
% nah
%  & Jahr
%  & 4 \\
% groß
%  & Koalition
%  & 4 \\
%\end{tabular}

		\bottomrule
	\end{tabularx}
	\caption{Kollokation Adjektive + Nomen, die häufigsten zehn Kombinationen.}
	\label{tab:adjektive_nomina}
\end{table}

\subsection{Häufigste Nomina}
Für die Bestimmung der häufigsten Nomina, die nicht gleich Ministerpräsident
sind, wurden folgende Befehle verwendet:
\begin{verbatim}
	> n = @[pos='N(E|N)' & word!='Ministerpräsident.*'];
	> group n target lemma cut 35 > "tabelle2.tex";
\end{verbatim}

Mit dem ersten Befehl werden alle Nomina (POS gleich NN oder NE), die nicht
gleich Ministerpräsident, beziehungsweise einer Deklination von
Ministerpräsident, sind abgefragt und in der Variable n gespeichert. Mit dem
zweiten Befehl wurden die so abgefragten Nomen nach ihrem Lemma grupiert und
nach ihrer Häufigkeit sortiert. Mit \it{cut 35} wurde die Ausgabe auf Nomen
beschränkt, die mindestens 35 mal auftreten. Somit wurde die Ausgabe auf die
häufigsten 10 Nomen beschränkt, wie gefordert. Der Wert von 35 wurde wie im
vorigen Beispiel manuell durch einen unbeschränkten Aufruf des selben
Befehls ermittelt. Die Ausgabe wurde in die Datei tabelle2.tex umgeleitet. 

Das Ergebnis ist der Tabelle \ref{tab:nomina} zu entnehmen.
% tabelle 2: nomina top10
\begin{table}[htpb]\label{t}
	\center
	\begin{tabularx}{0.35\textwidth}{lr}
		\toprule
		\textbf{Nomen} & \textbf{Quantität}\\
		\midrule
		% tabelle muss auf nur-Zeilen gekürzt werden!
		\begin{tabular}{llr}
(none) & Wulff
 & 112 \\
 & CDU
 & 87 \\
 & Christian
 & 65 \\
 & Land
 & 62 \\
 & Jahr
 & 62 \\
 & Landesbank
 & 54 \\
 & SPD
 & 48 \\
 & Deutschland
 & 40 \\
 & Euro
 & 39 \\
 & Regierung
 & 35 \\
\end{tabular}

		\bottomrule
	\end{tabularx}
	\caption{Nomina, die häufigsten zehn Vorkommen.}
	\label{tab:nomina}
\end{table}

\subsection{Adjektiv in Kollokation mit "`Ministerpräsident"'}
Für die Bestimmung der häufigsten Nomina, die nicht gleich Ministerpräsident
sind, wurden folgende Befehle verwendet:
\begin{verbatim}
	> adj = @[pos='ADJ(A|D)']@[word='Ministerpräsident.*'];
	> group adj match lemma cut 2 > "tabelle3.tex";
\end{verbatim}
% tabelle 3: nomina top10
Mit dem ersten Befehl werden alle Adjektive (POS gleich ADJA oder ADJD)
ausgelesen, die von "`Ministerpräsident"' oder einer Deklination des Wortes
gefolgtw erden. Mit dem zweiten Befehl werden die Ergebnisse nach dem Lemma des
Adjektivs gruppiert. Mit \it{cut 2} wird die Ausgabe auf Adjektive begrenzt, die
mindestens zweimal zusammen mit Ministerpräsident auftreten. Der Wert 2 wurde
wie in den vorigen Beispielen manuell ermittelt. Auch hier wurde die Ausgabe
wieder in eine Datei, diesmal tabelle3.tex, umgeleitet. 

Das Ergebnis ist der Tabelle \ref{tab:adj_mpraes} zu entnehmen.
\begin{table}[htpb]\label{t}
	\center
	\begin{tabularx}{0.5\textwidth}{lr}
		\toprule
		\textbf{Adjektiv} & \textbf{Quantität}\\
		\midrule
		% tabelle muss auf nur-Zeilen gekürzt werden!
		% Adjektive + Suchwort (McAllister)
% Die ersten 10 Zeilen der CQP Ausgabe
36-jährige
 & 2 \\

		\bottomrule
	\end{tabularx}
	\caption{Quantität der Adjektive, die in Kollokation mit dem Suchwort
	auftreten.}
	\label{tab:adj_mpraes}
\end{table}

\subsection{Subtitle}

Plain text.

\subsection{Another subtitle}

More plain text.

\end{document}

